\documentclass[12pt]{article}
\usepackage[utf8]{inputenc}
\usepackage[T1]{fontenc}
\usepackage{hyperref}
\title{Lokaa Language Resources for K-12 Curriculum Development}
\author{Compiled by Claude Sonnet 5.5 (claude-sonnet-5-5) for Andr\'as Kornai}
\date{}
\begin{document}
\maketitle

\section{Introduction}
Lokaa (ISO 639-3 yaz; Glottocode loka1252) is an Upper Cross River language of the Yakurr people of Cross River State, Nigeria. Search results also give it as Yak\"o, Lok\"o and L\`ok\'ee. Very little material exists. This report found no school materials and no recent materials of any kind.

\section{Primary Educational Resources}
No primers, textbooks or school programmes were found in the search results.

\section{Linguistic and Reference Materials}
\begin{itemize}
\item Alexander Ubi Iwara, \emph{Phonology and grammar of L\`ok\'ee: a preliminary study} (SOAS MPhil thesis, 1982). Its appendix has sample texts in a preliminary orthography. \url{https://eprints.soas.ac.uk/28582/}
\item Wikipedia article on the Yak\"o language, with phoneme inventories. It rests on dated sources. \url{https://en.wikipedia.org/wiki/Yak%C3%B6_language}
\item Glottolog bibliographic entry for Bendor-Samuel, ``Yakurr syllable patterns'' (1969). \url{https://glottolog.org/resource/reference/id/40951}
\end{itemize}
A 1974 chapter by Elaine Berry on Yakurr tone appears in \emph{Ten Nigerian Tone Systems}. The search results gave no URL for it.

\section{Religious and Community Resources}
None were found in the search results.

\section{Digital and Audio Resources}
None were found in the search results.

\section{Curriculum Development Context}
No dictionary, standard orthography or classroom material was found. A curriculum would need new primers, a settled orthography and a lexicon. The 1982 thesis is the only orthographic starting point found. Checking SIL, Nigerian language boards and Yakurr community bodies is a next step, but this search did not cover them.

\end{document}